\begin{figure}[ht]
\centering
\begin{tikzpicture}[
  vertex/.style = {circle, draw=engblack, thick, fill=englime!25,
                   inner sep=0pt, minimum size=7mm, font=\small},
  leaf/.style   = {circle, draw=engblue, very thick, fill=engblue!15,
                   inner sep=0pt, minimum size=7mm, font=\small},
  edge/.style   = {thick, draw=engblack},
]
% A labelled tree on six vertices and the first step of its Prufer encoding.
% The smallest-labelled leaf is vertex 1; its neighbour 4 is recorded.
\node[leaf]   (v1) at (0, 0)     {$1$};
\node[vertex] (v4) at (2, 0)     {$4$};
\node[vertex] (v3) at (4, 0.9)   {$3$};
\node[vertex] (v6) at (4, -0.9)  {$6$};
\node[leaf]   (v2) at (6, 1.6)   {$2$};
\node[leaf]   (v5) at (6, 0.2)   {$5$};

\draw[edge] (v1) -- (v4);
\draw[edge] (v4) -- (v3);
\draw[edge] (v4) -- (v6);
\draw[edge] (v3) -- (v2);
\draw[edge] (v3) -- (v5);

% Mark the leaves and the first recorded label.
\node[font=\scriptsize, anchor=north, engblue] at (v1.south) {leaf};
\node[font=\scriptsize, anchor=south, engblue] at (v2.north) {leaf};
\node[font=\scriptsize, anchor=north, engblue] at (v5.south) {leaf};
\node[font=\scriptsize, anchor=south west, engblue]
  at ($(v4.north east) + (0, 0.05)$) {record $4$};
\end{tikzpicture}
\caption{A labelled tree on six vertices and the first step of the
Prufer encoding. The smallest-labelled leaf is vertex $1$; the
algorithm removes it and records its lone neighbour $4$, then repeats
on the remaining tree. Six vertices give a length-four sequence over
the labels $\set{1, \ldots, 6}$, and the map between trees and
sequences is a bijection, so the tree count is $6^{4} = 6^{6-2}$,
Cayley's formula. The blue vertices are the current leaves; the first
recorded symbol is the neighbour of the smallest of them.}
\label{fig:vol02:ch17:prufer}
\end{figure}
